\section{Conclusion}
\label{sec:conclusion}

We presented an empirical evaluation of hybrid offline-to-online RL (H2O+) for bus holding control in a multi-fidelity setting: an offline dataset of past operations combined with a lower-fidelity simulator that supplies online interactions, evaluated against a higher-fidelity SUMO microsimulation.
The contribution is the recipe (which H2O+ components to use under which fidelity gap) and the evidence (where H2O+ is competitive with pure offline, pure online, and the RE-SAC SUMO expert checkpoint on transit holding under deployment constraints, where it is not, and what the operationally honest trade-offs are), rather than a proposal of new algorithmic modifications.

On a SUMO microsimulation built on the real topology of a 12-line, 389-bus corridor with a $675{,}000$-transition offline buffer, H2O+ reaches $-658$K under the deployment-relevant SIM-online setting (single SUMO seed, $5$ training seeds), narrowly above the RE-SAC SUMO expert checkpoint included in $\Doff$ (referred to as ep39 throughout, $-666$K) and $\approx 12\%$ above pure offline Cal-QL with full data ($-749$K) \emph{at the default operating point only} --- as we report below, this single-scenario ordering between H2O+ and ep39 does not survive the $3 \times 3$ stress grid.
Pure-online SAC trained from scratch fails at the $200$-epoch budget used for parity ($-1{,}654$K mean over $5$ seeds, with one seed diverging below the no-intervention baseline); a longer $1{,}000$-epoch retrain partially recovers but remains below H2O+ on the shaped composite reward (\Cref{tab:multiseed}).
On the multi-scenario stress test (3 SUMO seeds $\times$ 3 OD scales, paired CRN, \Cref{tab:multiseed}) the picture is more nuanced: H2O+, $1{,}000$-epoch pure-online SAC, and the RE-SAC SUMO expert ep39 form a near-tie within $\approx 100$K of one another, while H2O+ reliably beats the unlearned controllers (zero-hold, BC, Daganzo) by $25$--$33\%$.
H2O+ does \emph{not} consistently surpass ep39 across the $9$ scenarios tested; the ``slightly above ep39'' single-scenario number does not extrapolate.
Against five standard offline-to-online RL baselines (IQL, AWAC, TD3+BC, RLPD, WSRL) at three seeds each, all run with \texttt{snapshot\_reset} off to match the H2O+ Contrastive configuration, H2O+ clearly dominates the pure-offline IQL/AWAC ($+380$ to $+474$K paired mean), is statistically indistinguishable from RLPD on this benchmark, and is slightly outperformed by WSRL ($-15$K, $p = 0.021$) --- evidence that on transit holding the explicit dynamics-discriminator and IS reweighting do not provide an additional benefit beyond a hybrid offline+online buffer with standard SAC updates.
The operationally honest takeaway is that H2O+ produces a learned policy that recovers most of the RE-SAC SUMO expert's performance under a tunable shaped reward and reliably beats unlearned baselines, not that it consistently improves over that expert.
The ablation sweep ($\approx$30 configurations) surfaces three practical findings: (a) the optimal discriminator depends on the online--offline fidelity gap (Contrastive InfoNCE for SIM-online, TransitionDiscriminator for SUMO-online); (b) the offline-Q floor is a mixed-effect stabiliser, helping some configurations and hurting others; (c) the factored DynamicsDiscriminator from H2O+'s original prescription performs worst in either setting.
We deliberately avoid the term ``sim-to-real'' because the evaluation environment is itself a simulator; establishing true sim-to-real transfer requires historical-data replay against operational AVL/APC records, route-level travel-time calibration of $\Mreal$, weekend/incident-day scenarios, and an explicit calibration vs.\ test-day separation---all listed in \Cref{sec:discussion} as requirements for future validation.
Multi-seed evaluation across SUMO seeds and OD-scale stress tests, and passenger-side operational metrics (waiting time, headway CV, Jain fairness, large-gap rate), are reported in \Cref{tab:multiseed}; the calibration roadmap above is what remains.

Beyond the immediate transit application, the two structural issues isolated here---weak action-conditional signal under sparse actions, and reward-scale drift under multi-fidelity training---correspond to features that appear across multi-fidelity RL domains (ramp metering, signal phase selection, demand-responsive transit), so the diagnostic procedure and the recipe-by-fidelity-regime guidance may transfer to those settings.
Extending the evaluation to operational data, multi-fidelity simulator hierarchies, and joint timetable-and-holding optimisation are natural next directions.
